\section[Rédaction Bourbaki \redac{324} : rang constant en caractéristique $p$]{Rédaction Bourbaki \redac{324} : rang constant en caractéristique $p$ sur un corps imparfait\protect\piste\protect\verifie{Redaction324}}\label{sec:bourbaki-324}

La rédaction \redac{324}, \q{Variétés différentiables et analytiques. Fascicule de résultats (état 1)}, des archives Bourbaki (Archives Henri Poincaré, \ahp{755}) est de Serre et datée de juin 1959. Sa transcription reste locale pour l'instant ; on n'en cite ici aucun passage.

Le corps de base $K$ est un corps valué complet. Au §~4, n\textsuperscript{o}~9 (pp.~44--45 du fac-similé), le texte énonce en caractéristique nulle l'équivalence de deux conditions : $f$ est de rang $r$ au voisinage de $P$ ; $f$ s'écrit $z'_i = z_i$ ($i \le r$), $z'_i = 0$ ($i > r$) dans des coordonnées locales convenables. Dans ses notes finales (p.~94, note~5 du §~4), le rédacteur observe qu'en caractéristique $p$ la seconde condition entraîne la première mais non l'inverse, $x \mapsto x^p$ étant partout de rang nul. Il propose de caractériser plutôt les applications qui s'écrivent localement $z'_i = z_i^{p^{k(i)}}$ ($i \le r$), $z'_i = 0$ ($i > r$), et se demande si ce sont exactement celles qui sont localement de rang constant pour les dérivations de toutes hauteurs.

Sur un corps imparfait la réponse est non, pour toute notion de rang qui ne dépend que du germe de $f$ à translation près.

\begin{encadre}
\begin{contrexemple}[à la caractérisation envisagée dans la rédaction \redac{324}]\mbox{}
Soit $K$ un corps valué complet non discret de caractéristique~$2$ et $t \in K$ qui n'est pas un carré, par exemple $K = \mathbf F_2(\!(t)\!)$. L'application polynomiale $f(x, y) = x^2 + t y^2$ de $K^2$ dans $K$ est additive, donc $f(P + h) = f(P) + f(h)$ et $f$ a le même germe en tout point à translation près ; sa différentielle est partout nulle. Elle est injective. Elle n'admet pourtant, au voisinage d'aucun point, d'expression $z'_1 = z_1^{2^k}$ ou $z'_1 = 0$ dans des coordonnées locales.
\end{contrexemple}
\end{encadre}

\begin{proof}
En caractéristique~$2$, $(x + x')^2 = x^2 + x'^2$, d'où l'additivité, et $d f = 2x\,dx + 2ty\,dy = 0$. Si $x^2 + t y^2 = 0$ avec $y \ne 0$, alors $t = (x/y)^2$, exclu ; donc $f$ a un noyau nul, et elle est injective. Une expression de la forme proposée, la cible étant de dimension~$1$, fait dépendre $f$ de la seule coordonnée $z_1$ sur un ouvert $U$ : deux points de $U$ de même $z_1$ et de $z_2$ distincts, qui existent puisque l'image de $U$ par la carte est ouverte et que $K$ n'a pas de point isolé, ont même image, contre l'injectivité.
\end{proof}

Le même exemple vaut en caractéristique $p$ avec $x^p + t y^p$, $t \notin K^p$ : il repose sur le fait que $\{1, t\}$ est une $p$-base, de sorte que $f$ est même une bijection de $K^2$ sur $K$. Sur un corps parfait cet obstacle disparaît, puisque $x^p + c\,y^p = (x + c^{1/p} y)^p$ ; la question y reste ouverte ici.

\begin{sloppypar}
\noindent\emph{Vérification en Lean.} L'application est \texttt{sqForm} ; \texttt{sqForm\_add}, \texttt{hasFDerivAt\_sqForm} et \texttt{sqForm\_injective} donnent l'additivité, la différentielle nulle et l'injectivité, \texttt{not\_injOn\_of\_factor\_fst} l'argument de non-injectivité, \texttt{sqForm\_not\_normalForm} la conclusion, et \texttt{laurent\_X\_not\_square}, \texttt{sqForm\_laurent\_injective} l'exemple $\mathbf F_2(\!(X)\!)$, dans \href{\leanblobmain/Redaction324.lean}{\texttt{Redaction324.\allowbreak{}lean}}. Les cartes sont réduites à ce que l'argument utilise (une application d'image ouverte) ; ni les variétés analytiques ni les dérivations de Hasse--Schmidt ne sont formalisées.
\end{sloppypar}

\litterature{Serre, \emph{Lie Algebras and Lie Groups} (1964), partie~II, ch.~III, où le théorème du rang est énoncé en caractéristique nulle ; Bourbaki, \emph{Variétés différentielles et analytiques, fascicule de résultats} (1967), §~5 ; EGA~IV, §~17, et EGA~$0_{\mathrm{IV}}$, §~21 ($p$-bases) ; les travaux de Dieudonné sur les groupes formels en caractéristique $p$ (1955--1957). On n'a trouvé dans la littérature ni la question du rédacteur ni cet exemple dans le cadre des variétés analytiques.}
